\section{Implementation}
\label{sec:implementation}

This section describes how the proposed methodology is realised in software. The system
is implemented in Python and is exposed both as a set of command-line entry points and as
a browser-based operator console. The implementation follows the module boundaries
introduced in Section~\ref{sec:architecture}: generation, evaluation, execution,
orchestration, monitoring, and the user interface.

\subsection{Technology stack and repository layout}
\label{subsec:impl-stack}

The public side is defined with the AWS Cloud Development Kit for Python and synthesized
to CloudFormation; the private side is defined with Ansible playbooks. Static analysis
relies on established open-source scanners, cost estimation on a cost-analysis tool, and
live compliance on the cloud provider's configuration service. The operator console is
built with Streamlit. Table~\ref{tab:repo} lists the directories that implement each
concern.

\begin{table}[ht]
\centering
\caption{Repository layout and the concern each directory implements.}
\label{tab:repo}
\begin{tabular}{|L{5.5cm}|L{8.7cm}|}
\hline
\textbf{Directory} & \textbf{Contents} \\
\hline
\texttt{generation/} & Provider backends and the CDK regeneration loop
(\texttt{run\_cdk\_regen.py}). \\
\hline
\texttt{security\_gate/} & The gate (\texttt{iac\_security\_gate.py}) and the scanner adapters
under \texttt{security\_gate/scanners/}. \\
\hline
\texttt{execution/} & The subprocess execution helper
(\texttt{exec\_code.py}). \\
\hline
\texttt{pipeline/} & Orchestration, credentials, and operational-loop emitters. \\
\hline
\texttt{monitoring/} & CloudTrail, dashboards, the operational-loop function, and SSM
registration. \\
\hline
\texttt{scripts/} & \texttt{run\_cdk\_pipeline.py} and \texttt{run\_hybrid\_pipeline.py}. \\
\hline
\texttt{ui/} & The Streamlit console and the multi-project registry. \\
\hline
\texttt{risk\_scoring/} & Dataset preparation and model training. \\
\hline
\end{tabular}
\end{table}

Figure~\ref{fig:impl-modules} shows how these directories depend on one another at
runtime. The command-line entry points in \texttt{scripts/} sit at the top; they call
into \texttt{pipeline/} for orchestration, which in turn drives generation
(\texttt{generation/}), evaluation (\texttt{security\_gate/}), and monitoring (\texttt{monitoring/}). Every
directory that runs an external process funnels through the single execution helper in
\texttt{execution/}, and the operator console in \texttt{ui/} reuses the same entry
points rather than duplicating logic.

\begin{figure}[ht]
\centering
\begin{tikzpicture}[node distance=8mm and 10mm]
\node[proc] (ui) {\texttt{ui/}};
\node[proc, right=of ui] (scripts) {\texttt{scripts/}};
\node[proc, below=of scripts] (pipe) {\texttt{pipeline/}};
\node[proc, below left=of pipe] (aigen) {\texttt{generation/}};
\node[proc, below=14mm of pipe] (eval) {\texttt{security\_gate/}};
\node[proc, below right=of pipe] (mon) {\texttt{monitoring/}};
\node[proc, below=16mm of aigen] (exec) {\texttt{execution/}};
\draw[flow] (ui) -- (scripts);
\draw[flow] (scripts) -- (pipe);
\draw[flow] (pipe) -- (aigen);
\draw[flow] (pipe) -- (eval);
\draw[flow] (pipe) -- (mon);
\draw[flow] (aigen) -- (exec);
\draw[flow] (eval) -- (exec);
\draw[flow] (pipe.south) to[out=-90,in=90] (exec.north);
\end{tikzpicture}
\caption{Runtime module dependencies. Entry points drive orchestration, which drives
generation, evaluation, and monitoring; all external processes pass through the single
execution helper.}
\label{fig:impl-modules}
\end{figure}

\subsection{Generation layer}
\label{subsec:impl-generation}

Two interchangeable backends implement the generation contract behind a common
interface, so that the choice of provider is a configuration flag rather than a code
change. Both expose the same functions to call the model, extract the code from the
response, and save the result, and both centralise their default model identifiers while
allowing an explicit override. Provider-specific error handling (for example, quota and
throttling detection) is isolated behind the provider boundary, and credentials are read
from explicit arguments first and then from the environment. The regeneration loop
described in Algorithm~\ref{alg:regen} orchestrates these backends and persists the
per-attempt artifacts.

\subsection{Security gate and scanner adapters}
\label{subsec:impl-gate}

The gate is implemented in \texttt{iac\_security\_gate.py}, which collects the
synthesized templates, runs the heuristic checks and the scanner adapters, deduplicates
the findings, computes the score of Equation~\ref{eq:score}, applies the decision
function of Equation~\ref{eq:decision}, and writes the gate report. Each scanner lives in
its own adapter and returns a dictionary with at least a status and a message; no adapter
ever raises, so a missing tool degrades to a well-defined status while the gate proceeds.
Table~\ref{tab:adapters} lists the adapters.

\begin{table}[ht]
\centering
\caption{Scanner adapters and the signal each contributes.}
\label{tab:adapters}
\begin{tabular}{|L{6cm}|L{7.8cm}|}
\hline
\textbf{Adapter} & \textbf{Signal} \\
\hline
\texttt{checkov\_adapter.py} & Policy-as-code findings on the synthesized templates. \\
\hline
\texttt{cfn\_lint\_adapter.py} & CloudFormation best-practice and validity findings. \\
\hline
\texttt{infracost\_adapter.py} & Estimated monthly cost, used as the cost dimension. \\
\hline
\texttt{aws\_config\_adapter.py} & Count of non-compliant live-compliance rules. \\
\hline
\texttt{ansible\_lint\_adapter.py} & Findings on the on-premises playbooks. \\
\hline
\texttt{secret\_scan\_adapter.py} & Regular-expression secret detection (no external
tool). \\
\hline
\texttt{ml\_risk\_adapter.py} & Logistic-regression probability of insecurity, converted
to points. \\
\hline
\end{tabular}
\end{table}

An abbreviated gate report is shown in Listing~\ref{lst:report}. The
\texttt{components} object exposes the per-dimension contribution to the score, which
makes the decision explainable, and the \texttt{scanner\_status} object records which
tools ran.

\begin{lstlisting}[language=Python,caption={Abbreviated gate report.},label={lst:report}]
{
  "run_id": "cdk_20260614T093934Z",
  "timestamp": "2026-06-14T09:39:42+00:00",
  "score": 50,
  "decision": "review",
  "message": "Manual review required before deployment.",
  "components": {"severity": 30, "cost": 10, "aws_config": 10, "ml_risk": 0},
  "findings": [
    {"severity": "high", "source": "checkov", "category": "sg_ssh_open",
     "resource_id": "MySG", "template": "App.template.json",
     "message": "Security group allows ingress on port 22 from 0.0.0.0/0"}
  ],
  "scanner_status": {"checkov": "ok", "cfn_lint": "ok",
                     "infracost": "ok", "aws_config": "ok", "ml_risk": "ok"}
}
\end{lstlisting}

\subsection{Orchestration and governance}
\label{subsec:impl-orchestration}

Two command-line entry points drive the pipeline. The CDK-only entry point,
\texttt{run\_cdk\_pipeline.py}, orchestrates bootstrap, synthesis, gating, diff, and
deployment for the public side, and can optionally invoke the regeneration loop on a
reject. The hybrid entry point, \texttt{run\_hybrid\_pipeline.py}, runs the same
gate, risk-scoring, deployment, and observability workflow over both a bring-your-own CDK
directory and a bring-your-own Ansible directory --- with no code-generation step ---
gating each side independently with the same engine and writing a combined report. The
orchestration routes all external commands through the execution backbone and honours the
deployment-governance rules of Section~\ref{subsec:method-governance}. Listing
\ref{lst:cli} shows representative invocations.

\begin{lstlisting}[language=bash,caption={Representative pipeline invocations.},label={lst:cli}]
# CDK-only: synth, gate, and deploy if the gate allows
python scripts/run_cdk_pipeline.py --project-dir generated_cdk --deploy

# Hybrid: gate both sides; scope the Ansible scan to files changed since HEAD~1
python scripts/run_hybrid_pipeline.py \
    --cdk-path examples/hybrid-demo/cdk \
    --ansible-path examples/hybrid-demo/ansible \
    --base-ref HEAD~1

# Read-only status of every gated target
python scripts/run_hybrid_pipeline.py --query-status
\end{lstlisting}

On the hybrid side, the Ansible scan can be scoped to the files changed since a given git
reference, which limits analysis to the modified plays. The decision returned by each
entry point is surfaced as a distinct process return code, so that both the console and
any external continuous-integration system can react to \emph{pass}, \emph{review}, and
\emph{reject} without parsing free text.

Figure~\ref{fig:impl-sequence} traces the hybrid pipeline as a sequence of interactions
between the orchestrator, the gate, the cloud and on-premises targets, and the monitoring
layer, showing where the decision gates the deployment step on each side independently.

\begin{figure}[ht]
\centering
\begin{tikzpicture}[node distance=6mm and 8mm, font=\small]
\node[proc] (orch) {Orchestrator\\\texttt{run\_hybrid}};
\node[proc, right=20mm of orch] (gate) {Gate\\\texttt{security\_gate/}};
\node[proc, right=of gate] (cdk) {AWS\\(CDK)};
\node[proc, right=of cdk] (ans) {On-prem\\(Ansible)};
\node[proc, below=16mm of gate] (emit) {Emit\\(SSM/EventBridge/CW)};
\draw[flow] (orch) -- node[above,font=\scriptsize]{templates} (gate);
\draw[flowback] (gate) -- node[below,font=\scriptsize]{report} (orch);
\draw[flow] (orch.north) to[out=30,in=150] node[above,font=\scriptsize]{deploy if pass} (cdk.north);
\draw[flow] (orch.north) to[out=20,in=160] node[above,font=\scriptsize,pos=0.6]{deploy if pass} (ans.north);
\draw[flow] (gate) -- (emit);
\draw[flow] (cdk.south) |- (emit.east);
\draw[flow] (ans.south) |- ($(emit.east)+(0,-3mm)$);
\end{tikzpicture}
\caption{Hybrid pipeline interaction sequence. The orchestrator gates both sides with the
same engine, deploys each side only on an accepting decision, and emits telemetry.}
\label{fig:impl-sequence}
\end{figure}

\subsection{Execution backbone}
\label{subsec:impl-exec}

All external commands (CDK, Ansible, and the scanners) run through a single subprocess
helper in \texttt{exec\_code.py}. The helper merges standard output and standard error to
preserve the ordering of events, optionally streams output line by line for live
rendering in the console, and returns the stable \texttt{\{return\_code, output\}}
contract that every caller relies on. Concentrating process execution in one place keeps
the output-capture and error-propagation behaviour consistent across the whole system.

% \subsection{Operational-loop observability}
% \label{subsec:impl-monitor}

% The monitoring layer implements three complementary observability concerns. Gate metrics
% (score, decision, and the component breakdown) are published to a custom CloudWatch
% namespace; compliance drift is watched by the configuration service and reacted to by an
% event-driven function; and an audit trail of API activity is captured by CloudTrail. A
% dashboard presents the risk score and the decision per target for both the public and the
% private side, and the operational-loop function reacts to drift and rollback events to
% enable automated remediation. The event shapes and parameter layout emitted by the
% orchestration layer are consumed unchanged by this layer, which keeps the two decoupled.

\subsection{Operator console and multi-project registry}
\label{subsec:impl-ui}

The Streamlit console is a thin orchestration layer over the command-line entry points:
every privileged action is delegated to a subprocess with credentials injected through
the environment, never as command-line arguments. The console presents the workflow as a
sequence of tabs --- settings, generation, gate review, plan/diff, deploy, and monitor
--- and renders the gate report, the decision, and the live monitoring reads. A
multi-project registry allows several independent projects to coexist: the default
project keeps a flat log layout for backward compatibility, while any other project owns
a separate artifact directory so that gate reports, approvals, rejections, and monitoring
history never mix between projects. The active project's log directory is exported to the
subprocesses through an environment variable, which is how per-project isolation is
enforced end to end.
